
\chapter{素数定理的详细估计}
\label{app:pnt_estimates}

本附录补充第~\ref{ch:prime_number_theorem}~章证明提要中略去的技术细节：
$\zeta(1+it)\neq 0$ 的 $3$--$4$--$1$ 论证、零点自由区与围道边估计的要点，
以及 Tauberian 路线的简述。
\textbf{不}替代 Titchmarsh、Davenport、Edwards 等专著中的完整证明；
目标是使本书读者能回溯关键不等式与估计的出处与逻辑。

%===========================================================
\section{\texorpdfstring{$\zeta(1+it)\neq 0$}{zeta(1+it)≠0}：三角不等式路线}
\label{sec:app-no-zero-Re1}
%===========================================================

\subsection{Mertens 型非负三角式}
对实数 $\theta$，有
\begin{equation}
  3+4\cos\theta+\cos 2\theta
  =2(1+\cos\theta)^{2}
  \ge 0.
  \label{eq:app-trig}
\end{equation}
因而对 $\sigma>1$ 与实数 $t$，
\begin{equation}
  \sum_{n\ge 1}
  \frac{\Lambda(n)}{n^{\sigma}}
  \bigl(3+4\cos(t\log n)+\cos(2t\log n)\bigr)
  \ge 0.
  \label{eq:app-mertens-sum}
\end{equation}
左端可写为
\begin{equation}
  -3\frac{\zeta'(\sigma)}{\zeta(\sigma)}
  -4\operatorname{Re}\frac{\zeta'(\sigma+it)}{\zeta(\sigma+it)}
  -\operatorname{Re}\frac{\zeta'(\sigma+2it)}{\zeta(\sigma+2it)}.
  \label{eq:app-logderiv-combo}
\end{equation}

\subsection{化为 \texorpdfstring{$\zeta$}{zeta} 的模}
对 $\sigma>1$ 取实部并指数化，可得 Hadamard 形式的不等式
（常称为 $3$--$4$--$1$ 不等式）
\begin{equation}
  \zeta(\sigma)^{3}\,
  \bigl|\zeta(\sigma+it)\bigr|^{4}\,
  \bigl|\zeta(\sigma+2it)\bigr|
  \ge 1.
  \label{eq:app-341}
\end{equation}
（由~\eqref{eq:app-logderiv-combo} 积分或对 $\log|\zeta|$ 的相应组合得到；
完整推导见 Titchmarsh 第 3 章或 Davenport 第 13 章。）

\subsection{反证}
设 $t\neq 0$ 且 $\zeta(1+it)=0$。
因 $\zeta$ 在 $s=1$ 仅有单极点，在 $s=1+it$ 全纯，故零点阶 $m\ge 1$。
令 $\sigma\to 1^+$：
\begin{itemize}
  \item $\zeta(\sigma)\sim 1/(\sigma-1)$，故 $\zeta(\sigma)^{3}\asymp (\sigma-1)^{-3}$；
  \item $|\zeta(\sigma+it)|\asymp (\sigma-1)^{m}$；
  \item $|\zeta(\sigma+2it)|$ 趋于 $|\zeta(1+2it)|$（有限，因 $1+2it\neq 1$ 且由归纳/已知在 $\operatorname{Re}s=1$ 上至多有限阶零点的局部分析；
    严格处理见标准教材）。
\end{itemize}
于是~\eqref{eq:app-341} 左端
$\asymp (\sigma-1)^{-3+4m}\cdot O(1)$。
若 $m\ge 1$，则 $-3+4m\ge 1$，当 $\sigma\to 1^+$ 时左端 $\to 0$，与 $\ge 1$ 矛盾。
故 $\zeta(1+it)\neq 0$。

\begin{remark}{本书正文的用法}{}
第~\ref{ch:prime_number_theorem}~章只引用结论~\eqref{eq:app-341} 与反证结构；
数值常数与 $\zeta(\sigma+2it)$ 的细致极限需查 Titchmarsh §3.3--3.4。
\end{remark}

%===========================================================
\section{零点自由区与 \texorpdfstring{$\psi(x)=x+o(x)$}{psi=x+o(x)}}
\label{sec:app-zero-free}
%===========================================================

\subsection{de la Vallée Poussin 型自由区}
存在绝对常数 $c>0$，使得在区域
\begin{equation}
  \sigma\ge 1-\frac{c}{\log\bigl(|t|+2\bigr)}
  \label{eq:app-free-region}
\end{equation}
内 $\zeta(\sigma+it)\neq 0$（$|t|$ 充分大时）。
证明仍依赖对数导数的估计与 $3$--$4$--$1$ 型不等式的定量形式
（Titchmarsh 第 3 章；Davenport 第 13 章）。

\subsection{对显式公式截断的应用}
在自由区~\eqref{eq:app-free-region} 下，将 Perron 积分截断至高度 $T$ 并向左平移至
$\sigma=1-c'/\log T$，可得
\begin{equation}
  \psi(x)
  =x
  +O\bigl(x\,e^{-c''\sqrt{\log x}}\bigr)
  \label{eq:app-psi-error}
\end{equation}
（经典 de la Vallée Poussin 误差形；常数可改进）。
特别地 $\psi(x)=x+o(x)$，即素数定理。

第~\ref{ch:riemann_explicit_formula}~章在\textbf{全部}非平凡零点上写出精确公式；
此处仅需自由区内「没有零点」以保证截断误差 $o(x)$。

%===========================================================
\section{围道边估计（要点）}
\label{sec:app-contour}
%===========================================================

矩形围道（第~\ref{ch:integral_transforms}~章与附录~\ref{app:contour}）上：
\begin{enumerate}
  \item \textbf{水平边} $s=\sigma\pm iT$，$-1\le\sigma\le c$：
    需 $|\zeta'/\zeta(\sigma+iT)|\ll (\log T)^{O(1)}$（在避开零点的高度 $T$ 上）；
    与 $x^{\sigma}/|s|$ 结合使积分 $O(x^{c}/T\cdot\mathrm{poly}\log T)$ 等。
  \item \textbf{左竖直边}：在自由区左界上，$x^{\sigma}$ 给出 $x^{1-c/\log T}$ 型压制。
\end{enumerate}
选取 $T=T(x)\to\infty$ 的合适速率即得~\eqref{eq:app-psi-error}。
完整 $\varepsilon$--$T$ 管理见 Titchmarsh 第 3、7 章。

%===========================================================
\section{Tauberian 路线（简述）}
\label{sec:app-tauberian}
%===========================================================

Wiener--Ikehara 定理：若
\[
  F(s)=\sum_{n\ge 1}a_n n^{-s}
  \qquad (a_n\ge 0)
\]
在 $\operatorname{Re}s>1$ 收敛，且 $F(s)-A/(s-1)$ 可连续延拓至 $\operatorname{Re}s\ge 1$，
则 $\sum_{n\le x}a_n\sim A x$。

取 $a_n=\Lambda(n)$，$F=-\zeta'/\zeta$，$A=1$。
由定理~\ref{thm:no-zero-on-one} 与 $\zeta$ 在 $\operatorname{Re}s=1$、$s\neq 1$ 处全纯，
$-\zeta'/\zeta-1/(s-1)$ 在该直线上连续（局部），
故 $\psi(x)\sim x$。

Newman 给出更短的轮廓积分证明（同样以 $\zeta(1+it)\neq 0$ 为输入）。
第~\ref{ch:integral_transforms}~章提及 Laplace / Ikehara 作为背景；
\textbf{本书主线仍以 Perron + 零点自由区 + 显式公式为主}，
Tauberian 路线作为等价的现代简证备查。

%===========================================================
\section*{本附录与正文的分工}
%===========================================================

\begin{center}
\begin{tabular}{ll}
\hline
\textbf{内容} & \textbf{位置} \\
\hline
PNT 逻辑骨架、主项留数、$o(x)$ 结论 & 第~\ref{ch:prime_number_theorem}~章 \\
$3$--$4$--$1$ 不等式与反证细节 & 本附录 §\ref{sec:app-no-zero-Re1} \\
零点自由区与误差~\eqref{eq:app-psi-error} & 本附录 §\ref{sec:app-zero-free} \\
围道边估计 & 本附录 §\ref{sec:app-contour}；附录~\ref{app:contour} \\
Ikehara / Newman & 本附录 §\ref{sec:app-tauberian} \\
显式公式（主项$+$全部零点） & 第~\ref{ch:riemann_explicit_formula}~章 \\
\hline
\end{tabular}
\end{center}
